%%%%%%%%%%  scan idx 340  |  folio 327  %%%%%%%%%%
\SGAPrintPage{327}%
\sgaStableTarget{sga6.e05.eq.3-5-1}\begin{align}
 (1+Xt)\nabla(1+Yt)
 &=\rho\!\left(\rho^{-1}(1+Xt)\circ\rho^{-1}(1+Yt)\right) \notag\\
 &=\rho\!\left(\frac{1+(X+1)t}{1+t}\circ
                 \frac{1+(Y+1)t}{1+t}\right) \notag\\
 &=\rho\!\left(\frac{(1+(X+1)(Y+1)t)(1+t)}
                      {(1+(X+1)t)(1+(Y+1)t)}\right) \notag\\
 &=\frac{1+(X+Y+XY)t}{(1+Xt)(1+Yt)} . \tag{3.5.1}
\end{align}

If \(K\) is a \(\lambda\)-ring, then for all \(x,y\in K\) we have:
\sgaStableTarget{sga6.e05.eq.3-5-2}\begin{equation}\tag{3.5.2}
  \gamma_t(xy)=\gamma_t(x)\nabla\gamma_t(y).
\end{equation}

We denote by \(P'_n\) the \(n\)-th universal polynomial of the law \(\nabla\); it is a polynomial in \(\mathbb Z[X_i,Y_j]_{1\le i,j\le n}\) whose isobaric components have weights ranging from \(n\) to \(2n\) with respect to all the variables \(X_i,Y_j\), (1.15) and which is symmetric with respect to these two groups of variables. For \(x,y\in K\), we therefore have:
\sgaStableTarget{sga6.e05.eq.3-5-3}\begin{equation}\tag{3.5.3}
  n\in\mathbb N,\quad \gamma^n(xy)
  =P'_n(\gamma^1(x),\ldots,\gamma^n(x),
       \gamma^1(y),\ldots,\gamma^n(y)).
\end{equation}

Finally, note that the law \(\nabla\) has a unit element, namely \(\rho_t(1+t)\), that is, \(1/(1-t)\); it will sometimes be denoted by \(\underline{1}_\nabla\).

\sgaStableTarget{sga6.e05.para.3-6}3.6. The ``operations \(\lambda\)'' are defined by the commutative diagram:
\sgaStableTarget{sga6.e05.eq.3-6-1}\[
\begin{tikzcd}[column sep=tiny]
1+K[[t]]^+ \arrow[r,"\lambda_u"] \arrow[d,"\rho_t"'] &
1+\{1+K[[t]]^+\}[[u]]^+ \arrow[d,"\Ghat(\rho_t)"']\\
1+K[[t]]^+ \arrow[r,"\lambda'_u"] &
1+\{1+K[[t]]^+\}[[u]]^+ ,
\end{tikzcd}
\tag{3.6.1}
\]
i.e. by the composite morphism of functors:
\[
  \lambda' : \Ghat\xrightarrow{\rho_t^{-1}}\Ghat\xrightarrow{\lambda_u}\Ghat\circ\Ghat\xrightarrow{\Ghat(\rho_t)}\Ghat\circ\Ghat ,
\]
and will be denoted by \(\lambda'\). They are defined by their value on \(1+Xt\); we have:
%%%%%%%%%%  scan idx 341  |  folio 328  %%%%%%%%%%
\SGAPrintPage{328}%
\[
\begin{aligned}
\lambda'_u(1+Xt) &=\Ghat(\rho_t)(\lambda_u(\rho_t^{-1}(1+Xt)))\\
&=\Ghat(\rho_t)\left(\lambda_u\left(\frac{1+(X+1)t}{1+t}\right)\right)\\
&=\Ghat(\rho_t)\left(\frac{\underline{1}_o+\{1+(X+1)t\}u}{\underline{1}_o+\{1+t\}u}\right)\\
&=\underline{1}_\nabla+\frac{\{1+Xt\}u}{\underline{1}_\nabla+\underline{1}_\nabla.u}\\
&=\rho_u^{-1}(\underline{1}_\nabla+\{1+Xt\}u).
\end{aligned}
\]

Denoting by \(\gamma'\) the operations \(\gamma\) corresponding to the operations \(\lambda'\), we therefore have:
\sgaStableTarget{sga6.e05.eq.3-6-2}\begin{equation}\tag{3.6.2}
  \gamma'_u(1+Xt)=\underline{1}_\nabla+\{1+Xt\}u .
\end{equation}

We denote by \(Q'_{i,j}\) the coefficient of \(t^j\) in \(\gamma'^i(1+\ldots+X_nt^n+\ldots)\); the reader can verify that it is a polynomial in \(Z[X_1,\ldots,X_{ij}]\) whose isobaric components have weights ranging between \(j\) and \(ij\).

Let \(K\) be a \(\lambda\)-ring, \(x\in K\). From diagram (3.6.1), we deduce the relations:
\sgaStableTarget{sga6.e05.eq.3-6-3}\begin{equation}\tag{3.6.3}
  \forall i\in\mathbb N,\quad\gamma_t(\lambda^i(x))=\lambda'^i(\gamma_t(x))~;
  \qquad
  \gamma_t(\gamma^i(x))=\gamma'^i(\gamma_t(x)).
\end{equation}
\sgaStableTarget{sga6.e05.eq.3-6-4}\begin{equation}\tag{3.6.4}
  \gamma^j(\gamma^i(x))
  =Q'_{i,j}(\gamma^1(x),\ldots,\gamma^{ij}(x)).
\end{equation}

Finally, we observe that the pre-\(\lambda\)-ring structure just defined on \(\Ghat\) is a \(\lambda\)-ring structure, since it is obtained from the one defined in 2.3 by the isomorphism \(\rho\).

\sgaStableTarget{sga6.e05.para.3-7}3.7. To avoid confusion, the functor \(\Ghat\), endowed with the structure of a functor to \(\lambda\)-rings defined in §2 (resp. in §3), will be denoted by \(\Ghat_o\) (resp. \(\Gnabla\)). It follows from the preceding considerations that to specify a
%%%%%%%%%%  scan idx 342  |  folio 329  %%%%%%%%%%
\SGAPrintPage{329}%
\(\lambda\)-ring structure on a ring \(A\), it is equivalent to specify a ring homomorphism \(A\to\Ghat(A)_o\) such that diagram 3.7.1 is commutative, or a ring homomorphism \(A\to\Ghat(A)_\nabla\) such that diagram 3.7.2 is commutative; the first defines the operations \(\lambda\), the second the operations \(\gamma\).
\sgaStableTarget{sga6.e05.eq.3-7-1}\[
\begin{tikzcd}[column sep=tiny]
A \arrow[r,"\lambda_u"] \arrow[d,"\lambda_t"'] & \Ghat^u(A)_o \arrow[d,"\Ghat^u(\lambda_t)"]\\
\Ghat^t(A)_o \arrow[r,"\lambda_u"] & \Ghat^u(\Ghat^t(A)_o)_o
\end{tikzcd}
\tag{3.7.1}
\]
\sgaStableTarget{sga6.e05.eq.3-7-2}\[
\begin{tikzcd}[column sep=tiny]
A \arrow[r,"\gamma_u"] \arrow[d,"\gamma_t"'] & \Ghat^u(A)_\nabla \arrow[d,"\Ghat^u(\gamma_t)"]\\
\Ghat^t(A)_\nabla \arrow[r,"\gamma'_u"] & \Ghat^u(\Ghat^t(A)_\nabla)_\nabla.
\end{tikzcd}
\tag{3.7.2}
\]

\sgaStableTarget{sga6.e05.para.3-8}3.8. Let \(K\) be a binomial ring (2.7); from (2.6.1) we deduce, \(\forall x\in K\), the relations:
\sgaStableTarget{sga6.e05.eq.3-8-1}\begin{equation}\tag{3.8.1}
  \gamma_t(x)=(1/(1-t))^x
  =\exp(-x\Log(1-t))
\end{equation}

Let \(a\) be an element of \(K\); multiplication by \(\gamma_t(a)\) under the law \(\nabla\) is characterized by its value on \(1+Xt\), namely:
\[
  \gamma_t(a)\nabla(1+Xt)=\rho_t\left(\lambda_t(a)\circ\frac{1+(X+1)t}{1+t}\right)=\rho_t\left(\left(\frac{1+(X+1)t}{1+t}\right)^a\right)
\]
\sgaStableTarget{sga6.e05.eq.3-8-2}\begin{equation}\tag{3.8.2}
  \gamma_t(a)\nabla(1+Xt)
  =(1+Xt)^a
  =\exp(a\Log(1+Xt))
  =\sum_{n=0}^{\infty}\binom{a}{n}X^nt^n .
\end{equation}
In particular, we see that the \(n\)-th universal polynomial defining multiplication by \(\gamma_t(a)\) under the law \(\nabla\) is the polynomial \(R_{n,a}\) (2.9).

%%%%%%%%%%  scan idx 343  |  folio 330  %%%%%%%%%%
\SGAPrintPage{330}%
\sgaStableTarget{sga6.e05.def.3-9}\textbf{Definition 3.9.} \uline{Let} \(K\) \uline{be a} \(\lambda\)-\uline{ring}; \uline{by an} augmented \(K\)-\(\lambda\)-algebra \uline{we mean the data of a} \(K\)-\(\lambda\)-\uline{algebra} \(A\) \uline{and a} \(\lambda\)-\uline{homomorphism of} \(K\)-\uline{algebras} \(A\to K\); a homomorphism of augmented \(K\)-\(\lambda\)-algebras \uline{is a} \(\lambda\)-\uline{homomorphism compatible with the augmentations}.

\sgaStableTarget{sga6.e05.ex.3-9-1}\textbf{Example 3.9.1.}: Let \((X,A)\) be a commutatively locally ringed topos, and \(K^\bullet(X)\) the Grothendieck group of the category of sheaves of locally free \(A\)-modules of finite type on \(X\). We have seen that \(K^\bullet(X)\) can be endowed with a pre-\(\lambda\)-ring structure (2.2.b, since every direct summand of a free module of finite type is then locally free of finite type) and that \(H^0(X,\underline{\mathbb Z})\) is canonically endowed with a binomial ring structure. We define a \(\lambda\)-homomorphism from \(H^0(X,\underline{\mathbb Z})\) to \(K^\bullet(X)\) as follows: let \(a\) be a section of the constant sheaf \(\underline{\mathbb Z}\) on \(X\); it defines a partition of the final object of \(X\) into a sum of disjoint opens \(U_n\), the section \(a\) taking the value \(n\) on the open \(U_n\); let \(\underline{F}\) be the sheaf equal to \(A^n\) on \(U_n\) for \(n\ge0\), and to \(0\) on \(U_n\) for \(n\le0\), and let \(\underline{F}'\) be the sheaf equal to \(A^{-n}\) on \(U_n\) for \(n\le0\), and to \(0\) on \(U_n\) for \(n\ge0\); we associate to \(a\) the element \(F-F'\) of \(K^\bullet(X)\).

It is immediate to verify that this does give a \(\lambda\)-homomorphism from \(H^0(X,\underline{\mathbb Z})\) to \(K^\bullet(X)\); we can then define the augmentation from \(K^\bullet(X)\) to \(H^0(X,\underline{\mathbb Z})\) by assigning to every locally free sheaf of finite type its rank, which is a section of \(\underline{\mathbb Z}\); additivity of rank implies that this defines a function on \(K^\bullet(X)\) with values in \(H^0(X,\underline{\mathbb Z})\). We have thus endowed \(K^\bullet(X)\) with an augmented pre-\(\lambda\)-algebra structure. This is always the structure referred to below.

%%%%%%%%%%  scan idx 344  |  folio 331  %%%%%%%%%%
\SGAPrintPage{331}%
\sgaStableTarget{sga6.e05.para.3-10}3.10
Let \(A\) be an augmented \(K\)-\(\lambda\)-algebra. We define a ring filtration on \(A\) as follows: for every \(n\ge0\), we denote by \(\Fil^n A\) the ideal generated by the products:
\sgaStableTarget{sga6.e05.eq.3-10-1}\begin{equation}\tag{3.10.1}
  \gamma^{i_1}(x_1)\ldots\gamma^{i_k}(x_k),
  \qquad \text{with } i_1+\ldots+i_k\ge n,\ \text{and } \forall i,\ \epsilon(x_i)=0
\end{equation}
\((\epsilon:A\to K\) being the augmentation).

Suppose that \(A\) is generated as an abelian group by a family of elements \((x_\alpha)_{\alpha\in I}\); then \(\Fil^n A\) is generated as a \(K\)-module by the products:
\sgaStableTarget{sga6.e05.eq.3-10-2}\begin{equation}\tag{3.10.2}
  \gamma^{i_1}(x_{\alpha_1}-\epsilon(x_{\alpha_1}))\ldots
  \gamma^{i_k}(x_{\alpha_k}-\epsilon(x_{\alpha_k})),
  \qquad \text{with } i_1+\ldots+i_k\ge n,\ \text{and } \alpha_i\in I
\end{equation}
Indeed, if \(x\) is an element of \(A\) with zero augmentation and we write \(x=\sum_i x_{\alpha_i}\), we have \(\sum_i \epsilon(x_{\alpha_i})=0\), hence \(x=\sum_i x_{\alpha_i}-\epsilon(x_{\alpha_i})\); furthermore, if \(a\in A\), we can write \(a=(a-\epsilon(a))+\epsilon(a)=\gamma^1(a-\epsilon(a))+\epsilon(a)\); using relations (3.4.1), we therefore see that the elements of \(\Fil^n A\) are linear combinations with coefficients in \(K\) of the elements (3.10.2); since these belong to \(\Fil^n A\), we obtain the required result.

The filtration thus defined is a ring filtration; it is decreasing, and we have:
\[
  \Fil^0A=A~;
  \qquad \Fil^1A=\Ker\epsilon.
\]
Moreover, it is functorial in \(A\): every \(\lambda\)-homomorphism of augmented \(K\)-\(\lambda\)-algebras is compatible with the filtrations. When we regard an augmented \(\lambda\)-ring as endowed with a filtration, this will mean the filtration above, unless expressly stated otherwise.

